\section{Whose effect is the assistant being asked to estimate?}
\label{exec22:targets-scope}

A banking assistant translating an offer proposal must specify whose outcome
the offer is intended to improve. ``The average effect is positive'' leaves
open whether the average describes existing customers, a new acquisition
population, an equal mix of branches, or a precision-weighted statistical
summary. These choices can lead to different actions even when every group
effect is known. Vilfort's \emph{Robust Inference for Weighted Estimands}
provides a useful statistical treatment of this ambiguity
\citep{exec22_weighted_vilfort}. It is an econometric contribution to specifying
and checking targets, rather than an algorithm for language-model discovery.
Its value for a causal assistant begins after the group effects have a
defensible interpretation.

\subsection{Separate identification, aggregation and sampling uncertainty}

Let $\theta=(\theta_1,\ldots,\theta_K)'$ contain effects for $K$ customer groups,
with each entry measured in the same units and over the same horizon. A weight
vector $w$ defines
\begin{equation}
 \tau_w=w'\theta,\qquad \mathbf{1}'w=1.
 \label{exec22:targets-definition}
\end{equation}
Normalization alone permits negative entries. If the question concerns an
actual customer mixture, require the simplex
$\Delta_K=\{w:w\geq0,\mathbf{1}'w=1\}$. A signed regression contrast may be a
legitimate statistical target, but it should not be described as the mean
outcome of an unspecified population. Group definitions, treatment regimes,
calendar times and horizons must be compatible before aggregation. A weight
change cannot repair an effect whose identification assumptions have failed.

The exact experiment studied in the paper assumes
\begin{equation}
 \widehat\theta\sim N(\theta,\Sigma),\qquad \Sigma\succ0,
 \label{exec22:targets-experiment}
\end{equation}
where $K$ is fixed, $\Sigma$ is known, and the baseline weights and alternative
class are fixed and known. The full joint covariance matters; separate group
standard errors do not supply it. Let $\Lambda$ be a nonempty compact convex
class of alternative normalized weights. It represents stated ambiguity or
disagreement about the target, rather than an automatically identified set of
possible bank populations. A target outside this class receives no protection
from robustness calculated inside it. Choosing the class after inspecting
estimated effects also changes the statistical problem.

Consider an authored two-group illustration with stipulated effects
$\theta=(6,-2)'$ USD per customer and baseline
$w=(3/4,1/4)'$. Its contribution is $4$. The alternative population
$\lambda=(1/4,3/4)'$ has contribution $0$. At a stipulated cost of USD~1 per
customer, the respective net values are $3$ and $-1$. Nothing about this
disagreement requires a hallucination or an inaccurate group-effect estimate.
The populations answer different questions. These are exact project examples,
not observations from a bank or replications of the paper's applications
\citep{exec22_weighted_reference}.

One transparent ambiguity class moves a fraction $\epsilon$ of baseline
weight to any simplex distribution:
\begin{equation}
 \Lambda_\epsilon=(1-\epsilon)w+\epsilon\Delta_2,
 \qquad 0\leq\epsilon\leq1.
 \label{exec22:targets-class}
\end{equation}
Because the smallest and largest group contributions are $-2$ and $6$, its
effect and net-value ranges are
\begin{equation}
 \{\lambda'\theta:\lambda\in\Lambda_\epsilon\}
   =[4-6\epsilon,4+2\epsilon],\qquad
 \mathcal V_\epsilon=[3-6\epsilon,3+2\epsilon].
 \label{exec22:targets-value}
\end{equation}
The worst-case net value is strictly positive only for $\epsilon<1/2$.
The alternative population above first enters the class at $\epsilon=2/3$:
its first coordinate requires
$1/4\geq(1-\epsilon)3/4$. The break-even ambiguity level therefore has a clear
business meaning. It states how much customer-mix uncertainty the proposed
decision tolerates under stipulated effects. It does not estimate how uncertain
the true mix actually is.

Even a correctly calibrated interval for the baseline need not protect the
other population. Under the additional stipulated model
$\widehat\theta\sim N((6,-2)',I_2)$, the baseline estimator has variance $5/8$.
Its ordinary equal-tail 95\% interval is
$w'\widehat\theta\pm z_{.975}\sqrt{5/8}$. Coverage of the alternative target
$0$ is the analytic probability
\begin{equation}
 \Phi\!\left(z_{.975}-\frac{4}{\sqrt{5/8}}\right)
 -\Phi\!\left(-z_{.975}-\frac{4}{\sqrt{5/8}}\right)
 \simeq0.0009686482.
 \label{exec22:targets-wrong}
\end{equation}
Here $\Phi$ is the standard normal distribution function. Thus the interval
retains 95\% coverage for its declared baseline while covering this other
target about 0.0969\% of the time. This is an exact known-normal calculation,
not a simulation, a defect in baseline inference, or an execution of the
paper's robust procedure. It illustrates why an assistant must preserve the
population as part of the result.

\subsection{Use heterogeneity to understand weight sensitivity}

The paper separates effect heterogeneity from the distance between targets.
Define the covariance-metric heterogeneity and distance by
\begin{equation}
 H(\theta)=\left\{\min_\gamma
  (\theta-\mathbf{1}\gamma)'\Sigma^{-1}
  (\theta-\mathbf{1}\gamma)\right\}^{1/2},\qquad
 d_\Sigma(\lambda,w)=\{(\lambda-w)'\Sigma(\lambda-w)\}^{1/2}.
 \label{exec22:targets-geometry}
\end{equation}
The minimum centers the effects at their generalized least-squares constant.
Since $(\lambda-w)'\mathbf{1}=0$, subtracting any constant leaves the target
difference unchanged. Cauchy--Schwarz in the whitened coordinates gives the
paper's sharp bound
\begin{equation}
 |\tau_\lambda-\tau_w|
 \leq H(\theta)d_\Sigma(\lambda,w).
 \label{exec22:targets-bound}
\end{equation}
A banking interpretation follows directly: populations can disagree little
because effects are nearly homogeneous, because their weights are close in the
relevant metric, or both. None of these statements checks confounding.
In the example with $\Sigma=I_2$, the centering constant is $2$,
$H^2=32$, and $d_\Sigma^2=1/2$. The squared bound is $16$, exactly the squared
effect difference $(0-4)^2$. This attained witness checks the stated geometry.

Minimax-bias weights minimize the maximum covariance-metric distance to
$\Lambda$. For the two-group class at $\epsilon=2/3$, the first coordinate
ranges from $1/4$ to $11/12$. With identity covariance the center is
$7/12$, with effect $8/3$ and worst-case effect deviation $8/3$. This is the
center of the stipulated oracle range $[0,16/3]$. It illustrates the geometric
compromise, but minimizing this bias radius need not minimize total interval
length, and the compromise weights need not be the population the sponsor
wishes to serve. A precision or robustness preference is a decision input;
an assistant should expose it rather than silently substitute it for the
business target.

\subsection{State which confidence event is protected}

The paper constructs a one-sided upper confidence bound $U_\beta$ on
$H(\theta)$, with error budget $\beta$, from the generalized least-squares
residual statistic. Under the exact experiment its square has a noncentral
chi-square law with $K-1$ degrees of freedom and conventional noncentrality
$H(\theta)^2$. The paper's notation indexes this distribution by
heterogeneity itself. A software conversion must reconcile that notation with
the noncentrality accepted by a numerical library. Confusing $H$ and $H^2$
would change the bound. The source also specifies a zero-bound branch when
the observed statistic is sufficiently small; arbitrary root finding without
that branch would not implement the stated procedure.

Let $R_w=\max_{\lambda\in\Lambda}d_\Sigma(\lambda,w)$,
$B_\beta=U_\beta R_w$ and $s_w=\sqrt{w'\Sigma w}$. The paper's two-sided
interval uses the folded-normal critical value $c_{1-\alpha}(b)$, defined as
the $(1-\alpha)$ quantile of $|N(b,1)|$:
\begin{equation}
 C_w=\left[ w'\widehat\theta
       \mathbin{\pm}c_{1-\alpha}(B_\beta/s_w)s_w\right].
 \label{exec22:targets-ci}
\end{equation}
This combines sampling variation with a confidence bound on target
disagreement. The heterogeneity bound, class optimization and critical-value
calculation each require their own numerical checks.

Proposition~6 of the selected original establishes
\begin{equation}
 \Pr_\theta\{\lambda'\theta\in C_w\}\geq1-\alpha-\beta
 \quad\text{for every fixed }\lambda\in\Lambda\text{ and every }\theta.
 \label{exec22:targets-point}
\end{equation}
This uniform statement about each fixed target differs from the event that
one interval contains every target together. Appendix~C.1, Proposition~7,
gives the two-sided simultaneous guarantee
\begin{equation}
 \Pr_\theta\{\lambda'\theta\in C_w\ \text{for all }\lambda\in\Lambda\}
    \geq1-2\alpha-\beta.
 \label{exec22:targets-set}
\end{equation}
The one-sided simultaneous guarantee is $1-\alpha-\beta$. With
$\alpha=\beta=.05$, the stated lower guarantees are 90\% for each fixed
target and 85\% for the simultaneous two-sided event. They are lower bounds,
not assertions that actual coverage equals those numbers. To allocate a 5\%
budget to the latter event one possible choice is
$\alpha=.0125$, $\beta=.025$. Reporting ``95\% confidence'' without stating
the event and both budgets would conceal a substantive change in protection.

This distinction matters for agents that inspect many customer mixes before
selecting a favorable one. A guarantee for every preassigned target does not
automatically guarantee coverage of a target selected from the same results.
Simultaneous protection can protect such a selection within a fixed admissible
class, under its assumptions and its actual confidence budget. It does not
license learning the class from outcomes without further justification, nor
does it correct other forms of model or treatment selection.

\subsection{Move to feasible estimation without discarding the conditions}

The source develops asymptotic extensions with a uniform central limit theorem,
consistent covariance and weight estimates, eigenvalue bounds and regularity
for estimated classes. These conditions are material for dependent banking
panels, where covariance estimation is part of the scientific argument.
The main asymptotic development keeps $K$ fixed. A growing-$K$ appendix
requires separate technical assessment before use here. Treating a large
collection of sparse branch effects as though it automatically satisfied the
fixed-dimensional experiment would bypass that issue.

Estimated classes introduce another boundary. Consistency of a class does not
ensure that it contains every fixed target on the limiting boundary in a given
sample. The source's Proposition~U8 uses a containment condition; an enlarged
class with fixed positive slack is a sufficient route in its example.
Remark~8 explicitly distinguishes the theorem for any fixed $\delta>0$ from
applications reported at $\delta=0$. The observation that $\delta=.0001$
leaves displayed intervals unchanged is a numerical rounding observation,
rather than a proof of the theorem at zero slack. A proposed implementation
must state which class, containment argument and slack it actually uses.

For event studies, the required inputs are the group-time effect vector and
its full joint covariance, with compatible aggregation definitions. An interval
for each event time does not automatically become simultaneous protection
across all event times. The paper's public-data applications illustrate its
authors' results; without acquiring and inspecting their data, code and terms,
they do not establish a replication or a ready software bridge. A bounded
comparison should first use public known-law group effects, verify the oracle
geometry, exercise covariance and class-boundary failures, and then compare
the feasible procedures on matched targets.

The proposal should offer three explicit choices: commit to a substantively
justified fixed population; report a sensitivity range over declared weights;
or add a justified robust confidence construction when sampling uncertainty
also matters. The library, CLI and MCP representations must carry the group
definitions, weight role, full covariance, class, target-selection rule,
confidence event and budgets. A language model can help elicit these fields
and explain their consequences. It cannot infer missing population truth from
well-written prose or turn a valid interval for one target into evidence for
another.
